\documentclass[runningheads]{llncs}
\usepackage[T1]{fontenc}
\usepackage{graphicx}
\usepackage{amsmath,amssymb}
\usepackage{booktabs}

\begin{document}

\title{Title of Your Contribution\thanks{Supported by the example project.}}
\titlerunning{Short Title}   % shown in the page header when it is too long

\author{First Author\inst{1} \and
Second Author\inst{2}}
\authorrunning{F. Author and S. Author}

\institute{Example University, City, Country\\
\email{first.author@example.edu} \and
Example Lab, City, Country\\
\email{second.author@example.com}}

\maketitle

\begin{abstract}
Summarise the paper in 150 to 250 words: problem, method, result. Do not cite.
\keywords{First keyword \and Second keyword \and Third keyword.}
\end{abstract}

\section{Introduction}
Introduce the problem and state the contribution. Springer asks for numbered
references in order of appearance~\cite{knuth84}. Keep to the page limit of your
conference; references usually count unless the call for papers says otherwise.

\section{Preliminaries}
Definitions, theorems and proofs have their own environments in this class:

\begin{definition}
A \emph{widget} is a thing that is defined here once.
\end{definition}

\begin{theorem}\label{thm:main}
Every widget is a gadget.
\end{theorem}
\begin{proof}
Immediate from the definition.
\end{proof}

\section{Results}
Table~\ref{tab:res} shows the compact table style LNCS expects (captions above).

\begin{table}
  \caption{Example results.}\label{tab:res}
  \centering
  \begin{tabular}{lrr}
    \toprule
    Instance & Baseline & Ours \\
    \midrule
    small & 12 & 9 \\
    large & 340 & 210 \\
    \bottomrule
  \end{tabular}
\end{table}

\section{Conclusion}
What you showed, and what is left.

\bibliographystyle{splncs04}
\bibliography{refs}

\end{document}
